%! Function Model Selection and Construction — Unit 1, Lesson 1.7 — Student Packet Cover Sheet
\documentclass[10pt]{article}
\usepackage{precalculus-article}
\usepackage{precalculus-boxes}
\usepackage{ltablex}
\keepXColumns

\begin{document}

% Full-bleed plum banner
\begin{tikzpicture}[remember picture, overlay]
  \fill[plum] (current page.north west)
    rectangle ([yshift=-0.9in]current page.north east);
\end{tikzpicture}
\vspace{-0.8in}

\begin{center}
  {\color{white}\textbf{\LARGE Precalculus}}\\[4pt]
  {\color{lilacmid}\large\textbf{Unit 1: Functions and Change}}\\[6pt]
  {\color{white}\textbf{\large Lesson 1.7 \quad Function Model Selection and Construction}}
\end{center}

\vspace{0.22in}
\namedateperiod
\vspace{0.18in}

\begin{learningtargetbox}
\small
\begin{enumerate}[leftmargin=1.5em, itemsep=5pt]
  \item I can decide which quantity in a situation is the \textit{question} ---
        the model's input --- and which quantities become answers to it.
  \item I can choose a model type from a table, by taking first and second
        differences of equally spaced inputs.
  \item I can choose a model type from the context alone, by asking how the
        answer moves as the question grows.
  \item I can build a model by writing the constraint, solving it for the other
        quantities, and substituting into the one I want.
  \item I can use a model: evaluate it, find an average rate of change
        \textit{with units}, and state the domain the context allows and the
        assumption it rests on.
\end{enumerate}
\end{learningtargetbox}

\vspace{0.14in}

\begin{tocbox}
\small
\renewcommand{\arraystretch}{1.5}
\begin{tabularx}{\linewidth}{@{} c l X r @{}}
\rowcolor{lilac}
\textbf{\#} & \textbf{Component} & \textbf{Description} & \textbf{Score} \\
\midrule
1 & Warm-Up        & Spiral review of prerequisite skills                   & \blank{1.2cm} \\
2 & Guided Notes   & Picking, building, and using a function model          & \blank{1.2cm} \\
3 & Group Activity & The cardboard bin                                     & \blank{1.2cm} \\
4 & Exit Ticket    & Independent check on selection and construction        & \blank{1.2cm} \\
5 & Homework       & Practice and extension                                 & \blank{1.2cm} \\
\midrule
  & \multicolumn{2}{r}{\textbf{Total}} & \blank{1.2cm} \\
\end{tabularx}
\end{tocbox}

\vspace{0.14in}

\begin{remindbox}
\small
Every lesson so far handed you a function. Today you are handed a
\textbf{situation} --- 80 feet of fence, a barn wall, a pen --- and a situation
has several quantities in it at once. A function accepts only one. So the first
move in modeling is a decision: \textbf{which quantity is the question?}
Everything follows from that. The table's top row is somebody's answer to it;
the constraint is what forces the other quantities to answer to it; the domain
lists the questions the world actually allows; and the units on every rate name
the question you chose.
\end{remindbox}

\end{document}
